\documentclass[a4paper,11pt,oneside]{memoir}

% Castellano
\usepackage[spanish,es-tabla]{babel}
\selectlanguage{spanish}
\usepackage[utf8]{inputenc}
\usepackage{placeins}

\RequirePackage{booktabs}
\RequirePackage[table]{xcolor}
\RequirePackage{xtab}
\RequirePackage{multirow}

% Links
\usepackage[colorlinks]{hyperref}
\hypersetup{
	allcolors = {red}
}

% Ecuaciones
\usepackage{amsmath}

% Rutas de fichero / paquete
\newcommand{\ruta}[1]{{\sffamily #1}}

% Párrafos
\nonzeroparskip


% Imagenes
\usepackage{graphicx}
\newcommand{\imagen}[2]{
	\begin{figure}[!h]
		\centering
		\includegraphics[width=0.9\textwidth]{#1}
		\caption{#2}\label{fig:#1}
	\end{figure}
	\FloatBarrier
}

\newcommand{\imagenflotante}[2]{
	\begin{figure}%[!h]
		\centering
		\includegraphics[width=0.9\textwidth]{#1}
		\caption{#2}\label{fig:#1}
	\end{figure}
}



% El comando \figura nos permite insertar figuras comodamente, y utilizando
% siempre el mismo formato. Los parametros son:
% 1 -> Porcentaje del ancho de página que ocupará la figura (de 0 a 1)
% 2 --> Fichero de la imagen
% 3 --> Texto a pie de imagen
% 4 --> Etiqueta (label) para referencias
% 5 --> Opciones que queramos pasarle al \includegraphics
% 6 --> Opciones de posicionamiento a pasarle a \begin{figure}
\newcommand{\figuraConPosicion}[6]{%
  \setlength{\anchoFloat}{#1\textwidth}%
  \addtolength{\anchoFloat}{-4\fboxsep}%
  \setlength{\anchoFigura}{\anchoFloat}%
  \begin{figure}[#6]
    \begin{center}%
      \Ovalbox{%
        \begin{minipage}{\anchoFloat}%
          \begin{center}%
            \includegraphics[width=\anchoFigura,#5]{#2}%
            \caption{#3}%
            \label{#4}%
          \end{center}%
        \end{minipage}
      }%
    \end{center}%
  \end{figure}%
}

%
% Comando para incluir imágenes en formato apaisado (sin marco).
\newcommand{\figuraApaisadaSinMarco}[5]{%
  \begin{figure}%
    \begin{center}%
    \includegraphics[angle=90,height=#1\textheight,#5]{#2}%
    \caption{#3}%
    \label{#4}%
    \end{center}%
  \end{figure}%
}
% Para las tablas
\newcommand{\otoprule}{\midrule [\heavyrulewidth]}
%
% Nuevo comando para tablas pequeñas (menos de una página).
\newcommand{\tablaSmall}[5]{%
 \begin{table}
  \begin{center}
   \rowcolors {2}{gray!35}{}
   \begin{tabular}{#2}
    \toprule
    #4
    \otoprule
    #5
    \bottomrule
   \end{tabular}
   \caption{#1}
   \label{tabla:#3}
  \end{center}
 \end{table}
}

%
% Nuevo comando para tablas pequeñas (menos de una página).
\newcommand{\tablaSmallSinColores}[5]{%
 \begin{table}[H]
  \begin{center}
   \begin{tabular}{#2}
    \toprule
    #4
    \otoprule
    #5
    \bottomrule
   \end{tabular}
   \caption{#1}
   \label{tabla:#3}
  \end{center}
 \end{table}
}

\newcommand{\tablaApaisadaSmall}[5]{%
\begin{landscape}
  \begin{table}
   \begin{center}
    \rowcolors {2}{gray!35}{}
    \begin{tabular}{#2}
     \toprule
     #4
     \otoprule
     #5
     \bottomrule
    \end{tabular}
    \caption{#1}
    \label{tabla:#3}
   \end{center}
  \end{table}
\end{landscape}
}

%
% Nuevo comando para tablas grandes con cabecera y filas alternas coloreadas en gris.
\newcommand{\tabla}[6]{%
  \begin{center}
    \tablefirsthead{
      \toprule
      #5
      \otoprule
    }
    \tablehead{
      \multicolumn{#3}{l}{\small\sl continúa desde la página anterior}\\
      \toprule
      #5
      \otoprule
    }
    \tabletail{
      \hline
      \multicolumn{#3}{r}{\small\sl continúa en la página siguiente}\\
    }
    \tablelasttail{
      \hline
    }
    \bottomcaption{#1}
    \rowcolors {2}{gray!35}{}
    \begin{xtabular}{#2}
      #6
      \bottomrule
    \end{xtabular}
    \label{tabla:#4}
  \end{center}
}

%
% Nuevo comando para tablas grandes con cabecera.
\newcommand{\tablaSinColores}[6]{%
  \begin{center}
    \tablefirsthead{
      \toprule
      #5
      \otoprule
    }
    \tablehead{
      \multicolumn{#3}{l}{\small continúa desde la página anterior}\\
      \toprule
      #5
      \otoprule
    }
    \tabletail{
      \hline
      \multicolumn{#3}{r}{\small continúa en la página siguiente}\\
    }
    \tablelasttail{
      \hline
    }
    \bottomcaption{#1}
    \begin{xtabular}{#2}
      #6
      \bottomrule
    \end{xtabular}
    \label{tabla:#4}
  \end{center}
}

%
% Nuevo comando para tablas grandes sin cabecera.
\newcommand{\tablaSinCabecera}[5]{%
  \begin{center}
    \tablefirsthead{
      \toprule
    }
    \tablehead{
      \multicolumn{#3}{l}{\small\sl continúa desde la página anterior}\\
      \hline
    }
    \tabletail{
      \hline
      \multicolumn{#3}{r}{\small\sl continúa en la página siguiente}\\
    }
    \tablelasttail{
      \hline
    }
    \bottomcaption{#1}
  \begin{xtabular}{#2}
    #5
   \bottomrule
  \end{xtabular}
  \label{tabla:#4}
  \end{center}
}



\definecolor{cgoLight}{HTML}{EEEEEE}
\definecolor{cgoExtralight}{HTML}{FFFFFF}

%
% Nuevo comando para tablas grandes sin cabecera.
\newcommand{\tablaSinCabeceraConBandas}[5]{%
  \begin{center}
    \tablefirsthead{
      \toprule
    }
    \tablehead{
      \multicolumn{#3}{l}{\small continúa desde la página anterior}\\
      \hline
    }
    \tabletail{
      \hline
      \multicolumn{#3}{r}{\small continúa en la página siguiente}\\
    }
    \tablelasttail{
      \hline
    }
    \bottomcaption{#1}
    \rowcolors[]{1}{cgoExtralight}{cgoLight}

  \begin{xtabular}{#2}
    #5
   \bottomrule
  \end{xtabular}
  \label{tabla:#4}
  \end{center}
}




\graphicspath{ {./img/} }

% Capítulos
\chapterstyle{bianchi}
\newcommand{\capitulo}[2]{
	\setcounter{chapter}{#1}
	\setcounter{section}{0}
	\chapter*{#2}
	\addcontentsline{toc}{chapter}{#2}
	\markboth{#2}{#2}
}

% Apéndices
\renewcommand{\appendixname}{Apéndice}
\renewcommand*\cftappendixname{\appendixname}

\newcommand{\apendice}[1]{
	%\renewcommand{\thechapter}{A}
	\chapter{#1}
}

\renewcommand*\cftappendixname{\appendixname\ }

% Formato de portada
\makeatletter
\usepackage{xcolor}
\newcommand{\tutor}[1]{\def\@tutor{#1}}
\newcommand{\cotutor}[1]{\def\@cotutor{#1}}
\newcommand{\course}[1]{\def\@course{#1}}
\definecolor{cpardoBox}{HTML}{E6E6FF}
\def\maketitle{
  \null
  \thispagestyle{empty}
  % Cabecera ----------------
\noindent\includegraphics[width=\textwidth]{cabecera}\vspace{1cm}%
  \vfill
  % Título proyecto y escudo informática ----------------
  \colorbox{cpardoBox}{%
    \begin{minipage}{.8\textwidth}
      \vspace{.5cm}\Large
      \begin{center}
      \textbf{TFG del Grado en Ingeniería Informática}\vspace{.6cm}\\
      \textbf{\LARGE\@title{}}
      \end{center}
      \vspace{.2cm}
    \end{minipage}

  }%
  \hfill\begin{minipage}{.20\textwidth}
    \includegraphics[width=\textwidth]{escudoInfor}
  \end{minipage}
  \vfill
  % Datos de alumno, curso y tutores ------------------
  \begin{center}%
  {%
    \noindent\LARGE
    Presentado por \@author{}\\ 
    en Universidad de Burgos --- \@date{}\\
    Tutor: \@tutor{}\\
    Cotutor: \@cotutor{}\\
  }%
  \end{center}%
  \null
  \cleardoublepage
  }
\makeatother


% Datos de portada
\title{Investigación, creación y explotación de modelos para el estudio del COVID-19 en pacientes del Servicio Cántabro de Salud. \newline \newline Documentación Técnica}
\author{Segundo González González}
\tutor{Carlos López Nozal}
\cotutor{José Francisco Díez Pastor}
\date{\today}

\begin{document}

\maketitle



\cleardoublepage



%%%%%%%%%%%%%%%%%%%%%%%%%%%%%%%%%%%%%%%%%%%%%%%%%%%%%%%%%%%%%%%%%%%%%%%%%%%%%%%%%%%%%%%%



\frontmatter


\clearpage

% Indices
\tableofcontents

\clearpage

\listoffigures

\clearpage

\listoftables

\clearpage

\mainmatter

\appendix

\apendice{Plan de Proyecto Software}

\section{Introducción}

Como primera aproximación, en este apartado trataremos de explicar la planificación del proyecto y qué fases se han llevado a cabo.

\section{Planificación temporal}

Se ha considerado una metodología ágil (scrum) para llevar la planificación y el desarrollo de este proyecto, utilizando los conceptos teóricos y prácticos de la asignatura \textit{Gestión de proyectos}.

La plataforma elegida fue Github y en ella se fueron escribiendo tanto los \textit{milestones} como los \textit{issues}.

El proyecto se comenzó realizando sprints cada dos semanas y finalizando cada sprint con una reunión con los tutores, en las que se ponía en común los avances en el proyecto, proponiendo mejoras o cambios, y se establecían los objetivos para el siguiente sprint.

Estas sprints que comenzaron en septiembre de 2021 se vieron interrumpidos hasta junio de 2023 por diferentes motivos personales y se retomaron para poder terminar este proyecto.

A continuación, se detalla el trabajo realizado en cada sprint:

\subsection{14/09/2021 – 30/09/2021 \newline Sprint 0: Kickoff software project}

El primer sprint lo comenzamos con el número 0 debido a su carácter inicial que básicamente consistió en ver cómo se iban a planificar los sprints, exponer el alcance del proyecto y dar forma al inicio del proyecto de forma oficial.

\begin{itemize}
    \item Creación de repositorio en Github para poder llevar la planificación del proyecto, así como alojamiento de los futuros modelos desarrollados.
    \item Exposición del alcance del proyecto e información acerca de la cohorte de datos aún en proceso de petición al Instituto de de Investigación Sanitaria de Valdecilla (IDIVAL)
    \item Revisión preliminar de bibliografía en la que empezar a poner las bases de cómo acometer el proyecto.
    \item Consulta de TFG's presentados en la UBU relacionados con el proyecto a acometer  
    \item Formalización de registro del proyecto en UBU.
\end{itemize}


\subsection{01/10/2021 – 14/10/2021 \newline Sprint 1: Replicación trabajos relacionados}

Aquí se comienza a implantar los cimientos del proyecto comenzando una investigación profunda de trabajos relacionados que permitiese entender cómo acometer un proyecto de aprendizaje automático sobre datos sanitarios.

\begin{itemize}
    \item Investigación de bibliotecas en python para la interpretación de modelos de caja negra
    \item Escoger un trabajo relacionado con este proyecto para poder replicarlo
    \item Comenzar a añadir a la memoria los trabajos relacionados
\end{itemize}


\subsection{15/10/2021 – 31/10/2021 \newline Sprint 2: Cumplimentar primeros capítulos de la memoria}

Comenzamos a cumplimentar la memoria comenzando por el primer acercamiento a LaTeX y repasar la plantilla de la memoria.

\begin{itemize}
    \item Primeros pasos con LaTex y acercamiento a la utilización de la herramienta web Overleaf.
    \item Cumplimentar memoria: Introducción
    \item Cumplimentar memoria: Objetivos del proyecto
\end{itemize}


\subsection{01/11/2021 – 15/11/2021 \newline Sprint 3: Replica de trabajos relacionados}

Se recibe la cohorte de datos desde el IDIVAL y se comienza el desarrollo de un modelo de réplica del trabajo relacionado escogido \textit{Innate and Adaptive Immune Assessment at Admission to Predict Clinical Outcome in COVID-19 Patients}\cite{sansegundo:2021}

\begin{itemize}
    \item Replicar el modelo en python.
    \item Comparar resultados de la replicación con los obtenidos en los trabajos
    \item Comparación de resultados con otros clasificadores
\end{itemize}    


\subsection{16/11/2021 – 30/11/2021 \newline Sprint 4: Mejoras en la réplica del trabajo}

En este sprint se intenta mejorar el cuaderno jupyter de la réplica del trabajo, puesto que el resultado que arroja el actual cuaderno está bastante lejos del resultado del trabajo original.

\begin{itemize}
    \item Implementar Validación cruzada 2x5
    \item Utilizar pipeline en código
\end{itemize}


\subsection{01/12/2021 – 23/12/2021 \newline Sprint 5: Buscar clasificador con AutoSKlearn}

Se propone seguir con el concepto experimental del proyecto y se acude a la biblioteca autosklearn para buscar mejores clasificadores para nuestra cohorte de datos.

\begin{itemize}
    \item Estudiar uso de AutoSKlearn.
    \item Dockerización de Anaconda para incluir AutoSKLearn y demás librerías necesarias.
    \item Creación de cuaderno jupyter para uso de Autosklearn.
    \item Recoger resultados de AutoSKLearn aplicados a la cohorte de datos.
    \item Instalación y uso de librería Shap.
    \item Instalación y investigación de librerías de optimización de parámetros.
\end{itemize}


\subsection{21/06/2023 - 06/07/2023 \newline Sprint 6: Retomar el trabajo}

Se vuelve a retomar el trabajo dónde se dejó en 2022. Se revisan todos los cuadernos jupyter y la documentación para implementar los cambios que sean necesarios.

\begin{itemize}
    \item Revisión de la documentación de la memoria
    \item Reorganización del espacio GitHub
    \item Actualización contenedor con la última versión de Jupyter Notebooks
    \item Revisión y ejecución de cuadernos jupyter realizados
    \item Cumplimentación de los apartados de la memoria que aún no estaban documentados
    \item Realización de presentación para los vídeos a entregar
    \item Entrega de TFG en los espacios asignados en UBU Virtual
\end{itemize}


\section{Estudio de viabilidad}

A continuación se expone la posible viabilidad legal y económica de este proyecto.

\subsection{Viabilidad económica}

Por las características experimentales de este proyecto, creemos que este proyecto no se enmarca dentro de unas aplicaciones "comerciales"; el interés es puramente científico y en aras de aumentar el conocimiento sobre el aprendizaje automático con datos sanitarios.

También es correcto indicar, que llegado el momento en que se pudiera establecer un modelo de clasificación de pacientes en el triaje de una admisión hospitalaria con una muy alta precisión, es muy posible que las empresas de software hospitalario quisieran implementarlo en sus cores, si bien es cierto que tratándose en una pandemia, el uso del modelo distribuido debería ser de acceso universal.

El colectivo principal que mayor interés podría despertar en este proyecto sería, por supuesto, el de los profesionales sanitarios, sin dejar de lado a los entes gubernamentales, ya que estos últimos serían los posibles encargados de su distribución. 


\subsection{Viabilidad legal}

En el caso de las herramientas, no se considera ningún impedimento la legislación, puesto que todas las herramientas pueden utilizarse en cualquier aplicación comercial, debido a sus licencias (BSD o MIT en general).

Por contra, al estar tratando datos clínicos que se enmarcan dentro del nivel más alto de segurida de la información, se deberían establecer fuertes medidas de seguridad sobre ellos.
\apendice{Especificación de Requisitos}

\section{Introducción}

En este apartado se exponen los principales requisitos del proyecto.

\section{Objetivos generales}

Estos son los objetivos que se marcaron en el proyecto:

\begin{itemize}
    \item Replicar el trabajo \textit{Innate and Adaptive Immune Assessment at Admission to Predict Clinical Outcome in COVID-19 Patients}~\cite{sansegundo:2021}.
    \item Buscar trabajos relacionados y clasificar las variables, técnicas y herramientas que se utilizaron.
	\item Estudiar los datos proporcionados por el Instituto de Investigación Sanitaria de Valdecilla (IDIVAL) y considerar como procesarlos para poder obtener un un modelo de predicción que mejore los rendimientos del trabajo \textit{Innate and Adaptive Immune Assessment at Admission to Predict Clinical Outcome in COVID-19 Patients}~\cite{sansegundo:2021}.
    \item Utilizar librerías de aprendizaje automático que ayuden a encontrar el mejor método de clasificación de entre diferentes métodos a partir de las variables/marcadores para el Covid-19 y disponibles en el conjunto de datos.
    \item Minimizar las características a utilizar para su utilización en el modelo a desarrollar.
    \item Investigar qué hiperparámetros pueden mejorar el método de clasificación elegido para el modelo final.
\end{itemize}

\section{Catálogo de requisitos}

Debido al carácter fundamentalmente experimental de nuestro proyecto, distinguiremos en requisitos funcionales (qué realizan los modelos desarrollados) y no funcionales (cómo lo consigue).

\subsection{Requisitos funcionales}
\begin{itemize}

    \item \textbf{RF-1} Los datos deben poder ser divididos en conjuntos de entrenamiento y validación.
        \begin{itemize}
            \item \textbf{RF-1.1} Se indicará qué porcentaje de datos del total se utilizará en el conjunto de validación.
            \item \textbf{RF-1.2} Aleatorización de la inclusión de instancias en el conjunto de entrenamiento o validación, permitiendo siempre una división paritaria en ambos conjuntos.        
        \end{itemize}
    
    \item \textbf{RF-2} Se debe poder aplicar técnicas y algoritmos de preparación de los datos.
        \begin{itemize}
            \item \textbf{RF-2.1} Se debe poder eliminar los valores nulos de la característica a entrenar.
            \item \textbf{RF-2.2} Se debe poder imputar los valores perdidos o nulos en el resto del conjunto de datos
        \end{itemize}
    
    \item \textbf{RF-3} Se debe poder entrenar un modelo.
        \begin{itemize}
            \item \textbf{RF-3.1} Ofrecer la posibilidad de escoger entre diferentes algoritmos de clasificación.
            \item \textbf{RF-3.2} Debe guardarse el mejor modelo generado durante el entrenamiento.
            \item \textbf{RF-3.3} Se debe poder mostrar gráficamente los resultados del entrenamiento.
        \end{itemize}

    \item \textbf{RF-4} Se debe poder realizar búsquedas de mejor clasificador de forma automática.
        \begin{itemize}
            \item \textbf{RF-4.1} Se debe poder cambiar el dataset a utilizar.
            \item \textbf{RF-4.2} Se debe poder cambiar los parámetros de búsqueda.
            \item \textbf{RF-4.3} Se debe obtener una solución gráfica en la que observar los resultados.
        \end{itemize}

    \item \textbf{RF-5} Se debe poder optimizar los hiperparámetros de un modelo.
        \begin{itemize}
            \item \textbf{RF-5.1} Se debe poder cambiar el modelo a optimizar.
            \item \textbf{RF-5.2} Se debe obtener un gráfico con los resultados.
        \end{itemize}

    \item \textbf{RF-6} Se debe poder probar un modelo guardado.
        \begin{itemize}
            \item \textbf{RF-6.1} Se debe poder encontrar el modelo a cargar.
        \end{itemize}
\end{itemize}

\subsection{Requisitos no funcionales}

\begin{itemize}
    \item \textbf{RNF-1} Portabilidad: Se puede ejecutar entre diferentes sistemas operativos y entornos de desarrollo, lo que evita problemas de configuración y compatibilidad.
    \item \textbf{RNF-2} Reproducibilidad: El contenedor contiene todas las dependencias y configuraciones específicas, lo que facilita la reproducción de resultados y la colaboración.
    \item \textbf{RNF-2} Aislamiento: Los paquetes o bibliotecas instalados en el contenedor no afectarán a tu sistema host, y viceversa.
    \item \textbf{RNF-2} Escalabilidad: Permite manejar grandes conjuntos de datos y ejecutar cálculos intensivos de manera eficiente utilizando recursos distribuidos.
    \item \textbf{RNF-2} Administración simplificada: Gestión de múltiples Jupyter Notebooks y entornos. Se pueen crear y administrar diferentes contenedores para proyectos específicos o diferentes configuraciones de entorno.
    \item \textbf{RNF-2} Flexibilidad: Permite personalizar y configurar el entorno para instalar las bibliotecas y herramientas necesarias, configurar extensiones y temas, y ajustar los recursos del contenedor para optimizar el rendimiento.
\end{itemize}
\newpage

\section{Especificación de requisitos}

\subsection{Diagrama de caso de uso}

\imagen{diaguso}{Diagrama de caso de uso}

\subsection{Especificaciones de los casos de uso}

\tablaSinCabeceraConBandas{Caso de Uso \textit{División de datos en conjuntos de entrenamiento y validación}}{p{3cm} p{9cm}}{2}{casouso1}
{Caso de uso & División de datos en entrenamiento y validación \\ 
Requisitos &  RF-1\newline
            RF-1.1\newline
            RF-1.2 \\ 
Descripción & Los datos deben poder ser divididos en conjuntos de entrenamiento y validación. \\ 
Precondiciones & Se debe disponer de un conjunto de datos. \\ 
Acciones & 1. Se indicará qué porcentaje de datos del total se utilizará en el conjunto de validación. \newline
            2. Aleatorización de la inclusión de instancias en el conjunto de entrenamiento o validación, permitiendo siempre una división paritaria en ambos conjuntos. \\ 
Postcondiciones & Datos divididos en dos conjuntos.\\ 
Excepciones & El conjunto no lleva datos o está corrupto. \\ 
Importancia & Media\\
}

\tablaSinCabeceraConBandas{Caso de Uso \textit{Aplicar técnicas y algoritmos de preparación de los datos}}{p{3cm} p{9cm}}{2}{casouso2}
{Caso de uso & Preparación de datos \\ 
Requisitos &  RF-2\newline
            RF-2.1\newline
            RF-2.2 \\ 
Descripción & Los datos deben poder ser divididos en conjuntos de entrenamiento y validación. \\ 
Precondiciones & Se debe disponer de un conjunto de datos. \\ 
Acciones & 1.  Se debe poder eliminar los valores nulos de la característica a entrenar. \newline 
            2. Se debe poder imputar los valores perdidos o nulos en el resto del conjunto de datos. \\ 
Postcondiciones & Ddatos preparados.\\ 
Excepciones & El conjunto no lleva datos o está corrupto. \\ 
Importancia & Media\\
}

\tablaSinCabeceraConBandas{Caso de Uso \textit{Entrenar un modelo}}{p{3cm} p{9cm}}{2}{casouso3}
{Caso de uso & Entrenar un modelo \\ 
Requisitos &  RF-3\newline
            RF-3.1\newline
            RF-3.2\newline
            RF-3.3 \\ 
Descripción & Se debe poder entrenar un modelo. \\ 
Precondiciones & Se debe disponer de un conjunto de datos. \\ 
Acciones & 1. Ofrecer la posibilidad de escoger entre diferentes algoritmos de clasificación. \newline
           2. Debe guardarse el mejor modelo generado durante el entrenamiento. \newline
           3. Se debe poder mostrar gráficamente los resultados del entrenamiento. \\ 
Postcondiciones & Modelo entrenado.\\ 
Excepciones & El conjunto no lleva datos o está corrupto. \\ 
Importancia & Alta\\
}

\tablaSinCabeceraConBandas{Caso de Uso \textit{Búsqueda de mejor clasificador de forma automática.}}{p{3cm} p{9cm}}{2}{casouso4}
{Caso de uso & Búsqueda automática de mejor clasificador \\ 
Requisitos &  RF-4\newline
            RF-4.1\newline
            RF-4.2\newline
            RF-4.3 \\ 
Descripción & Se debe poder realizar búsquedas de mejor clasificador de forma automática. \\ 
Precondiciones & Se debe instalar la librería de automl. \\ 
Acciones & 1. Se debe poder cambiar el dataset a utilizar. \newline
           2. Se debe poder cambiar los parámetros de búsqueda. \newline
           3. Se debe obtener una solución gráfica en la que observar los resultados. \\ 
Postcondiciones & Mejor clasificador encontrado.\\ 
Excepciones & Error al realizar la tarea de búsqueda. \\ 
Importancia & Media\\
}

\tablaSinCabeceraConBandas{Caso de Uso \textit{Optimización de hiperparámetros.}}{p{3cm} p{9cm}}{2}{casouso5}
{Caso de uso & Optimización de hiperparámetros \\ 
Requisitos &  RF-5\newline
            RF-5.1\newline
            RF-5.2 \\ 
Descripción & Se debe poder optimizar los hiperparámetros de un modelo. \\ 
Precondiciones & Se debe instalar la librería de optimización de hiperparámetros. \\ 
Acciones & 1. Se debe poder cambiar el modelo a optimizar. \newline
           2. Se debe obtener un gráfico con los resultados. \\ 
Postcondiciones & Mejores hiperparámetros encontrados.\\ 
Excepciones & Error al realizar la tarea de optimización. \\ 
Importancia & Media\\
}

\tablaSinCabeceraConBandas{Caso de Uso \textit{Prueba de modelos}}{p{3cm} p{9cm}}{2}{casouso6}
{Caso de uso & Prueba de modelos \\ 
Requisitos &  RF-6\newline
            RF-6.1 \\ 
Descripción & Se debe poder probar un modelo guardado. \\ 
Precondiciones & Modelos desarrollados alojados correctamente \\ 
Acciones & 1. Se debe poder encontrar el modelo a cargar. \\ 
Postcondiciones & Prueba correcta de modelo\\ 
Excepciones & No se encuentra el modelo en la ubicación \\ 
Importancia & Media\\
}
\apendice{Especificación de diseño}

\section{Introducción}

En este apartado recogemos el diseño experimental desarrollado en el proyecto.

\section{Diseño de datos}

La cohorte de datos han sido proporcionados por el Instituto de Investigación Sanitaria de Valdecilla (IDIVAL) sobre una población de 305 pacientes, cuyas muestras fueron recogidas entre Abril de 2020 y Marzo de 2021 y que tuvieron un resultado positivo en COVID-19. El valor de clasificación es el \textbf{score} de evolución clínica que tiene un valor de 1 a 5, siendo 1=asintomático; 2=tratamiento con gafas nasales; 3=tratamiento con ventimask; 4=Ingreso en UCI; 5=exitus.

El dataset se compone de un archivo csv con las siguientes características:
\begin{itemize}
	\item 420 filas o determinaciones.
	\item 97 columnas o marcadores.
    \item Sólo hay 158 filas con el valor de clasificación score cumplimentado.
\end{itemize}

Los valores correspondientes a la intensidad de luz presentan una amplia variedad, moviéndose generalmente entre los intervalos de $-10^5$ a $10^5$. 

\imagen{ejemdatos}{Ejemplo de conjunto de datos}

\section{Diseño procedimental}

En la realización de un nuevo experimento, el científico de datos posee total libertad para elegir las herramientas y las técnicas que va a aplicar. En cualquier caso, suele coincidir de forma general el flujo de trabajo.

El proceso de creación de un experimento lleva habitualmente las siguentes fases: carga de datos, preparación de estos, configuración del modelo y entrenamiento. Puede verse en la imagen \nameref{fig:diagflujo1}.

\imagen{diagflujo1}{Diagrama de flujo de realización de un experimento}

\section{Diseño arquitectónico}
\apendice{Documentación técnica de programación}

\section{Introducción}

Comentaremos aquí los detalles más técnicos sobre el proyecto de forma que sea posible para cualquiera recrear los experimentos y continuar el trabajo aquí presentado.

\section{Estructura de directorios}

El código fuente del proyecto se encuentra en Github en el directorio \href{https://github.com/sgg0008/ml-COVID-19/tree/main/2. Cuadernos TFG}{2. Cuadernos TFG}

\section{Compilación, instalación y ejecución del proyecto}

El entorno de trabajo se  contenedor Docker \cite{docker:2013}. Bajo la aplicación \textit{Container Station} se instaló la imagen del contenedor \textit{jupyter/tensorflow-notebook} \cite{jupyterdocker:1991} que se descargó de la plataforma Docker Hub. En la instalación se añadió la creación de un volumen de almacenamiento persistente para poder ir guardando de forma segura los cuadernos que se iban realizando. Este contenedor se escogió porque ya incorpora gran parte del ecosistema de bibliotecas científicas de python; otras bibliotecas fueron instaladas en el arranque del contenedor ya que son utilizadas en este trabajo, como son Autosklearn, Pipeline profiler, Shap y Optuna.

\section{Manual del programador}

El entorno de ejecución es el mismo que el de desarrollo, realizándose en Jupyter Notebooks \cite{jupyter:2014} y en el lenguaje de programación Python \cite{python:1991}
\apendice{Documentación de usuario}

Debido al carácter plenamente experimental de este proyecto, no existe una distribución formal de los modelos desarrollados para un usuario convencional. Es por esto que este apartado no aplica en este proyecto.



\bibliographystyle{plain}
\bibliography{bibliografiaAnexos}

\end{document}
